\documentclass[11pt]{article}
\usepackage[margin=1in]{geometry}
\usepackage{amsmath,amssymb,amsthm}
\usepackage{hyperref}

\newtheorem{theorem}{Theorem}
\newtheorem{lemma}{Lemma}
\newtheorem{corollary}{Corollary}
\newtheorem{remark}{Remark}

\newcommand{\F}{\mathcal F}
\newcommand{\Lip}{\operatorname{Lip}}
\newcommand{\supp}{\operatorname{supp}}
\newcommand{\dist}{\operatorname{dist}}
\newcommand{\U}{\mathbb U}
\newcommand{\C}{\mathbb C}
\newcommand{\N}{\mathbb N}
\newcommand{\eps}{\varepsilon}

\title{Compact Lipschitz operators on unbounded spaces have only the periodic root spectrum}
\author{Run fa\_banach\_001, lane 4}
\date{2026-06-29}

\begin{document}
\maketitle

\section*{Source Question}

The source paper is Abbar, Coine, and Petitjean,
\emph{A note on the spectrum of Lipschitz operators and composition operators
on Lipschitz spaces}, arXiv:2310.08965.  Its first final open question asks
whether the bounded-space spectral corollary remains true for unbounded metric
spaces:

\begin{quote}
If \(\widehat f:\F(M)\to\F(M)\) is compact, is the set \(A\) of periods of
periodic points of \(f\) finite and
\[
\sigma(\widehat f)\setminus\{0\}=\bigcup_{n\in A}\U_n?
\]
\end{quote}
Here, as in the source paper's Corollary 3.1 and Theorem 2.5 setup, \(A\)
records periods of nonzero periodic points; the distinguished base point
\(0\) does not yield a nonzero Dirac eigenvector.

\section*{Definitions and Source Facts}

All metric spaces are pointed, complete, and have base point \(0\).  If
\(f:M\to M\) is base-point-preserving and Lipschitz, \(\widehat f\) denotes
the induced bounded operator on the Lipschitz-free space \(\F(M)\), defined on
Dirac masses by \(\widehat f(\delta(x))=\delta(f(x))\).  We write
\[
\U_n=\{z\in\C:z^n=1\}.
\]

We use four facts from the source paper and from Abbar--Coine--Petitjean,
\emph{Compact and weakly compact Lipschitz operators}, arXiv:2110.03231.

\begin{enumerate}
\item If \(x\) is a periodic point of exact period \(n\), then
\(\U_n\subset \sigma_p(\widehat f)\).
\item If \(T\) is compact, then
\(\sigma(T)\setminus\{0\}=\sigma_p(T)\setminus\{0\}\), and \(0\) is the only
possible accumulation point of \(\sigma(T)\).
\item If \(\widehat f\) is compact, then \(f\) satisfies the compactness
characterization \(P_1,P_2,P_3\): images of bounded subsets are totally
bounded, \(f\) is uniformly locally flat, and the tail condition \(P_3\) holds.
In particular, \(f\) is radially flat:
\[
 \frac{d(f(x),0)}{d(x,0)}\longrightarrow 0
 \qquad (d(x,0)\to\infty).
\]
\item The source paper uses cutoff operators \(W_R=\theta_R\widehat{\mathrm{Id}}\),
where \(\theta_R\) is \(0\) on \(B(0,R)\), linear on the annulus
\(R\le d(x,0)\le 2R\), and \(1\) outside \(B(0,2R)\).  These satisfy
\(\sup_R\|W_R\|\le 4\), \(W_R\gamma\) is supported outside \(B(0,R)\), and
\((I-W_R)\gamma\) is supported in \(B(0,2R)\).
\end{enumerate}

\section*{Proof Intuition}

The source proof for flat-at-infinity maps cuts an alleged unbounded
eigenvector into a far-out tail and tests the eigen-equation against a
functional supported on that tail.  Flatness at infinity makes the Lipschitz
constant of \(h\circ f\) on the tail arbitrarily small, forcing the nonzero
eigenvalue to be arbitrarily small.

Compactness gives less than flatness at infinity, but it gives exactly enough.
The \(P_3\) tail condition says that escaping domain pairs with image
separation comparable to domain separation must have image pairs with an
accumulation point.  Therefore, once at least one image remains far out, such a
large quotient is impossible.  This conditional tail-flatness is all the
cutoff proof needs, because the testing function vanishes on the bounded part
of the image.  Nonzero eigenvectors are forced into bounded supports; after
that, a large invariant closed ball reduces the problem to the bounded
corollary already proved in the source paper.

\section*{Tail Estimate}

\begin{lemma}[escaping-image tail flatness]\label{lem:tail}
Let \(\widehat f:\F(M)\to\F(M)\) be compact.  For every \(\eps>0\) there is
\(R_\eps>0\) such that whenever \(x,y\in M\) satisfy
\[
d(x,0)\ge R_\eps,\qquad d(y,0)\ge R_\eps,\qquad
\max\{d(f(x),0),d(f(y),0)\}\ge R_\eps,
\]
one has
\[
d(f(x),f(y))\le \eps\, d(x,y).
\]
\end{lemma}

\begin{proof}
If not, there is \(\eps_0>0\) and, for each \(k\), points \(x_k,y_k\) with
\(d(x_k,0),d(y_k,0)\ge k\), with at least one of \(d(f(x_k),0)\) and
\(d(f(y_k),0)\) at least \(k\), and with
\[
d(f(x_k),f(y_k))>\eps_0 d(x_k,y_k).
\]
The domain pairs escape to infinity.  The image pairs cannot have an
accumulation point in \(M\times M\), since one image coordinate escapes to
infinity.  This contradicts condition \(P_3\) in the compactness
characterization of arXiv:2110.03231, which would force a subsequence with
\(d(f(x_k),f(y_k))/d(x_k,y_k)\to0\).
\end{proof}

\begin{lemma}\label{lem:bounded-support}
If \(\widehat f\) is compact and
\(\widehat f(\gamma)=\lambda\gamma\) for some \(\gamma\ne0\) and
\(\lambda\ne0\), then \(\supp(\gamma)\) is bounded.
\end{lemma}

\begin{proof}
Assume, toward a contradiction, that \(\supp(\gamma)\) is unbounded.  Fix
\(\eps>0\).  By radial flatness and Lemma~\ref{lem:tail}, choose \(R\) so
large that
\[
d(f(x),0)\le \eps d(x,0)\quad(d(x,0)\ge R),
\]
\[
f(B(0,2R))\subset B(0,R),
\]
and the conclusion of Lemma~\ref{lem:tail} holds with \(R_\eps=R\).

Let \(W_R=\theta_R\widehat{\mathrm{Id}}\) be the source cutoff operator.
Since \(\supp(\gamma)\) is unbounded, \(W_R\gamma\ne0\).  Put
\[
C_R=M\setminus B(0,R).
\]
Choose \(g_R\in \Lip_0(\{0\}\cup C_R)\) with \(\|g_R\|_L=1\) and
\[
\langle g_R,W_R\gamma\rangle=\|W_R\gamma\|.
\]
Let \(h_R=W_R^*g_R=\theta_R g_R\).  Then \(h_R\in\Lip_0(M)\),
\(\|h_R\|_L\le4\), and \(h_R\) is supported outside \(B(0,R)\).

The element \((I-W_R)\gamma\) is supported in \(B(0,2R)\), hence
\(\widehat f((I-W_R)\gamma)\) is supported in \(B(0,R)\), where \(h_R\)
vanishes.  Therefore
\[
|\lambda|\,\|W_R\gamma\|
 = |\langle h_R,\widehat f(\gamma)\rangle|
 = |\langle h_R,\widehat f(W_R\gamma)\rangle|.
\]

It remains to estimate \(h_R\circ f\) on \(\{0\}\cup C_R\).  If
\(x,y\in C_R\) and both \(f(x),f(y)\in B(0,R)\), then
\(h_R(f(x))=h_R(f(y))=0\).  Otherwise at least one of \(f(x),f(y)\) lies
outside \(B(0,R)\), and Lemma~\ref{lem:tail} gives
\[
|h_R(f(x))-h_R(f(y))|
\le 4\,d(f(x),f(y))
\le 4\eps\,d(x,y).
\]
For the pair \(0,x\) with \(x\in C_R\), either \(h_R(f(x))=0\), or radial
flatness gives
\[
|h_R(f(x))-h_R(f(0))|
\le 4\,d(f(x),0)
\le 4\eps\,d(x,0).
\]
Thus
\[
\|(h_R\circ f)|_{\{0\}\cup C_R}\|_L\le 4\eps.
\]
Since \(W_R\gamma\in\F(\{0\}\cup C_R)\), we get
\[
|\lambda|\,\|W_R\gamma\|
 \le 4\eps\,\|W_R\gamma\|.
\]
Dividing by \(\|W_R\gamma\|\) gives \(|\lambda|\le4\eps\).  As \(\eps>0\)
was arbitrary, \(\lambda=0\), a contradiction.
\end{proof}

\section*{Main Result}

\begin{theorem}\label{thm:main}
Let \(M\) be a complete pointed metric space and let \(f:M\to M\) be a
base-point-preserving Lipschitz map.  If
\(\widehat f:\F(M)\to\F(M)\) is compact and \(A\) is the set of all natural
numbers \(n\) for which \(f\) has a nonzero periodic point of exact period
\(n\), then \(A\) is finite and
\[
\sigma(\widehat f)\setminus\{0\}=\bigcup_{n\in A}\U_n.
\]
Equivalently, since \(C_f=\widehat f^*\),
\[
\sigma(C_f)\setminus\{0\}=\sigma(\widehat f)\setminus\{0\}
=\bigcup_{n\in A}\U_n.
\]
\end{theorem}

\begin{proof}
The inclusion \(\supset\) is exactly the periodic-point eigenvector lemma from
the source paper.

Conversely, let \(\lambda\in\sigma(\widehat f)\setminus\{0\}\).  Since
\(\widehat f\) is compact, \(\lambda\in\sigma_p(\widehat f)\), so there is
\(\gamma\ne0\) such that \(\widehat f(\gamma)=\lambda\gamma\).  By
Lemma~\ref{lem:bounded-support}, \(\supp(\gamma)\) is bounded.

Choose \(R\) so large that \(\supp(\gamma)\subset B(0,R)\) and
\[
f(\overline{B}(0,R))\subset \overline{B}(0,R).
\]
This is possible by radial flatness.  Let \(M_R=\overline{B}(0,R)\), a
complete bounded pointed metric space, and consider \(f_R=f|_{M_R}:M_R\to
M_R\).  By the compactness characterization, \(f(M_R)\) is totally bounded and
\(f_R\) is uniformly locally flat.  Hence the source paper's bounded/totally
bounded spectral corollary applies to \(\widehat{f_R}:\F(M_R)\to\F(M_R)\).

The element \(\gamma\) belongs canonically to \(\F(M_R)\), and
\(\widehat{f_R}\gamma=\lambda\gamma\).  Therefore the source corollary gives
\(\lambda\in\U_n\) for some exact period \(n\) of a periodic point of \(f_R\),
hence of \(f\).  This proves
\[
\sigma(\widehat f)\setminus\{0\}\subset \bigcup_{n\in A}\U_n.
\]

Finally, \(A\) must be finite.  If infinitely many periods occurred, the
periodic-point lemma would put the infinite set \(\bigcup_{n\in A}\U_n\) into
\(\sigma_p(\widehat f)\).  This infinite subset of the unit circle has an
accumulation point different from \(0\), contradicting compactness of
\(\widehat f\).  The equality for \(C_f\) follows from
\(C_f=\widehat f^*\) and equality of spectra for an operator and its adjoint.
\end{proof}

\section*{Verification and Novelty Notes}

The proof uses only the source paper's support/cutoff and periodic-point
arguments plus the compactness characterization from arXiv:2110.03231.  The
new step is Lemma~\ref{lem:tail}, which extracts from \(P_3\) exactly the
tail estimate needed by the source's flat-at-infinity cutoff proof.

Cheap run indexes were searched for arXiv:2310.08965, the title keywords,
``compact Lipschitz operator spectrum'', ``roots of unity'', and related
Lipschitz-free keywords; no duplicate packet was found.  Web searches on
2026-06-29 for exact and close phrases around the final question found the
source arXiv record but no later answer.  Human review should check the
cutoff support step in Lemma~\ref{lem:bounded-support}, especially the use of
the source cutoff operator \(W_R\) and the passage to the bounded restriction.

\begin{thebibliography}{9}

\bibitem{ACP-spectrum}
A. Abbar, C. Coine, and C. Petitjean,
\emph{A note on the spectrum of Lipschitz operators and composition operators
on Lipschitz spaces}, arXiv:2310.08965.

\bibitem{ACP-compact}
A. Abbar, C. Coine, and C. Petitjean,
\emph{Compact and weakly compact Lipschitz operators}, arXiv:2110.03231.

\end{thebibliography}

\end{document}
